\documentclass[10pt]{article}
\usepackage[margin=0.68in]{geometry}
\usepackage{amsmath,amssymb,amsthm,mathtools}
\usepackage{microtype}
\AtBeginDocument{%
  \setlength{\abovedisplayskip}{6pt plus 2pt minus 2pt}%
  \setlength{\belowdisplayskip}{6pt plus 2pt minus 2pt}%
  \setlength{\abovedisplayshortskip}{4pt plus 2pt minus 2pt}%
  \setlength{\belowdisplayshortskip}{4pt plus 2pt minus 2pt}%
}

\newtheorem{theorem}{Theorem}
\newcommand{\N}{\mathbb N}
\newcommand{\dd}{\,d}
\newcommand{\floor}[1]{\lfloor #1\rfloor}
\newcommand{\fp}[1]{\{#1\}}

\title{A Floor-Saving Lower Bound for Unique Positive Differences\\
\large Proof sketch}
\author{GPT-5.5 Pro\thanks{Large language model by OpenAI; assisted with derivation and drafting through prompted interaction.}
\and
Lech Mazur\thanks{Human contributor, curator, editor, and responsible contact for public posting.}}
\date{}

\begin{document}
\maketitle

All logarithms are natural, and $\N=\{1,2,3,\ldots\}$.  Let
\[
        \lambda=\frac1{2\log2},\qquad
        \gamma=
        \lim_{N\to\infty}\left(\sum_{m\le N}\frac1m-\log N\right),
        \qquad
        B_1=\frac12+\gamma-\log(\log2).
\]

\begin{theorem}
Let $A\subseteq\N$ be such that, for every $n\ge1$, there is exactly one
ordered pair $(a,b)\in A^2$ with $a>b$ and $n=a-b$.  Write this pair as
$(a_n,b_n)$.  For every fixed $B>B_1$, arbitrarily large $n$ satisfy
\[
        a_n>\lambda\frac{n^2}{\log n-\log\log n+B}.
\]
Consequently,
\[
        \limsup_{n\to\infty}
        \frac{a_n(\log n-\log\log n+B_1)}{n^2}\ge\lambda .
\]
\end{theorem}

\begin{proof}[Sketch]
For a representation $n=t-b$, call $t$ the \emph{top}.  The spine is
\[
        \text{spacing}\longrightarrow \text{top count}\longrightarrow
        \text{integer floor}\longrightarrow \text{constant comparison}.
\]
The continuous majorant has the correct main term $X$.  The proof wins at the
next scale: the integer floor subtracts a deterministic fractional-part mass of
order $X/\log X$.

\smallskip
\noindent\textit{1. Contradiction setup and inverse weight.}
Fix $B>B_1$ and suppose that eventually
\[
        a_n\le \Phi(n),
        \qquad
        \Phi(x)=\lambda\frac{x^2}{H(\log x)},
        \qquad
        H(w)=w-\log w+B.                                      \tag{1}
\]
Choose $N_0\ge2$ so that (1) holds for $n\ge N_0$ and $\Phi$ is increasing on
$[N_0,\infty)$.  Put $T_*:=\Phi(N_0)$, let $\psi$ be the inverse of this branch,
and set $f(t)=1/\psi(t)$ for $t\ge T_*$, extending $f$ constantly on $[0,T_*]$.
Then $f$ is positive and nonincreasing.  Inverting
$t=\lambda r^2/H(\log r)$ gives
\[
        f(t)=\sqrt{\frac{2\lambda}{t\log t}}(1+o(1)),
        \qquad
        I(T):=\int_0^T f(t)\dd t
        =2\sqrt{\frac{2\lambda T}{\log T}}
        +o\!\left(\sqrt{\frac T{\log T}}\right).              \tag{2}
\]
All $o$-terms below are for $X\to\infty$, with $A$ and $B$ fixed.

\smallskip
\noindent\textit{2. Spacing.}
Choose an integer
\[
        Z_0>\max\left(T_*,\max_{1\le d<N_0}a_d\right).
\]
If $y<z$ are elements of $A$ and $z>Z_0$, then $(z,y)$ represents $d=z-y$.
Uniqueness forces $z=a_d$.  The choice of $Z_0$ excludes $d<N_0$, so (1) gives
$z\le\Phi(d)$ and hence $d\ge\psi(z)$.  Since $f$ is nonincreasing,
\[
        \int_y^z f(s)\dd s\ge (z-y)f(z)=d/\psi(z)\ge1.         \tag{3}
\]
Thus every large gap in $A$ costs one unit of $f$-mass.  Summing the gaps gives
\[
        |A\cap[1,T]|\le I(T)+O(1).                            \tag{4}
\]
We also need the weighted form.  If $z_1<\cdots<z_L$ are elements of $A$ above
$Z_0$ and $W\ge0$ is nonincreasing, then
\[
        \sum_{j=2}^L W(z_j)
        \le \int_{z_1}^{z_L} f(t)W(t)\dd t.                   \tag{5}
\]
Indeed, apply (3) to each gap $[z_{j-1},z_j]$ and use
$W(t)\ge W(z_j)$ on that gap.

\smallskip
\noindent\textit{3. Top count and floor.}
Let $X$ be a large integer, put $Y=\log X$ and $M=\Phi(X)$.  For all large
$X$, every difference $1\le n\le X$ has top at most $M$: use (1) for
$n\ge N_0$; the finitely many $n<N_0$ have fixed tops, eventually below $M$.
Hence
\[
        X=\sum_{\substack{t\in A:\,t\le M}} r_t,
        \qquad
        r_t:=\#\{b\in A:t-X\le b<t\}.                         \tag{6}
\]
This identity is exact: a counted pair gives a difference in $[1,X]$, and
uniqueness prevents double counting.  The terms with $t\le X$ are exactly
$\binom{|A\cap[1,X]|}{2}$, so (4) gives
\[
        \sum_{\substack{t\in A:\,t\le X}} r_t
        \le \frac12 I(X)^2+o(X/Y).                            \tag{7}
\]

For $t>X$ define
\[
        G_X(t)=\int_{t-X}^t f(s)\dd s .
\]
Since $f$ is nonincreasing, $G_X$ is nonincreasing.  If $t\in A$ and
$t>X+Z_0$, list the bottoms counted by $r_t$ as $b_1<\cdots<b_{r_t}$ and put
$b_{r_t+1}=t$.  Each upper endpoint in the gaps $[b_i,b_{i+1}]$ is above
$Z_0$, and these gaps lie in $[t-X,t]$.  Hence (3) gives
\[
        r_t\le\sum_{i=1}^{r_t}\int_{b_i}^{b_{i+1}} f(s)\dd s
        \le G_X(t).
\]
Since $r_t$ is an integer,
\[
        r_t\le \floor{G_X(t)}.                                \tag{8}
\]
This is the only non-continuous step.  Apply (5) to the large tops with
$W(t)=\floor{G_X(t)}$.  The first large top has at most $I(X)+O(1)$ bottoms,
and the fixed strip $X<t\le X+Z_0$ has $O(1)$ possible tops, each with
$O(I(X))$ bottoms.  By (2), these losses are $o(X/Y)$.  Therefore
\[
        X\le \frac12I(X)^2+
        \int_X^M f(t)\floor{G_X(t)}\dd t+o(X/Y).               \tag{9}
\]
Equivalently,
\[
        X\le J_B(X)-E(X)+o(X/Y),                               \tag{10}
\]
where
\[
        J_B(X)=\frac12I(X)^2+\int_X^M f(t)G_X(t)\dd t,
        \qquad
        E(X)=\int_X^M f(t)\fp{G_X(t)}\dd t .
\]

\smallskip
\noindent\textit{4. Continuous count.}
The continuous majorant is
\[
        J_B(X)=X+
        \lambda(3-2B-2\log(\log2))\frac{X}{Y}+o(X/Y).          \tag{11}
\]
Here is the constant ledger.  First,
\[
        \frac12I(X)^2=4\lambda X/Y+o(X/Y).                     \tag{12}
\]
Next split
\[
        \int_X^M f(t)G_X(t)\dd t
        =X\int_X^M f(t)^2\dd t+R(X),
\]
where
\[
        R(X)=\int_X^M f(t)\int_{t-X}^t(f(s)-f(t))\dd s\dd t .
\]
For the $f^2$ term use $t=\Phi(e^w)$ and write $w_0=\log\psi(X)$.  Then
\[
        w_0=Y/2+(\log Y)/2-(\log(2\lambda))/2+o(1),
        \qquad M=\Phi(e^Y).
\]
Also $f(\Phi(e^w))=e^{-w}$ and $dt=\Phi'(e^w)e^w\,dw$, so
\[
        f(t)^2\dd t
        =\lambda\left(\frac2{H(w)}-\frac{H'(w)}{H(w)^2}\right)\dd w.
\]
Thus
\[
        \int_X^M f(t)^2\dd t
        =\lambda\left[2\int_{w_0}^{Y}\frac{\dd w}{H(w)}
        +\frac1{H(Y)}-\frac1{H(w_0)}\right].                  \tag{13}
\]
Writing $w_0=Y/2+\rho_X$, the three $1/Y$ contributions inside the bracket are
\[
\begin{aligned}
        2\int_{w_0}^{Y}\frac{\dd w}{w}
        &=2\log2-\frac{4\rho_X}{Y}+o(1/Y),\\
        2\int_{w_0}^{Y}\frac{\log w-B}{w^2}\dd w
        &=\frac{2\log Y-4\log2+2-2B}{Y}+o(1/Y),\\
        \frac1{H(Y)}-\frac1{H(w_0)}
        &=-\frac1Y+o(1/Y).
\end{aligned}                                                  \tag{14}
\]
Since $\rho_X=(\log Y-\log(2\lambda))/2+o(1)$, the logarithms cancel and,
using $2\lambda\log2=1$,
\[
        X\int_X^M f(t)^2\dd t
        =X+\lambda(1-2B+2\log(\lambda/2))\frac{X}{Y}+o(X/Y).  \tag{15}
\]
For $R(X)$, put $t=Xu$ and $s=Xv$.  The local form of (2) on bounded
$u$-ranges, followed by truncation, gives
\[
        R(X)=\frac{2\lambda X}{Y}
        \int_1^\infty\left(\int_{u-1}^u\frac{\dd v}{\sqrt{uv}}-
        \frac1u\right)\dd u+o(X/Y)
        =2\lambda(2\log2-1)X/Y+o(X/Y).                       \tag{16}
\]
Adding (12), (15), and (16), and using
$4\log2+2\log(\lambda/2)=-2\log(\log2)$, gives (11).

\smallskip
\noindent\textit{5. Floor saving.}
The floor subtracts
\[
        E(X)=2\lambda(1-\gamma)X/Y+o(X/Y).                    \tag{17}
\]
Set
\[
        q(t)=Xf(t)=X/\psi(t).
\]
The range $\Phi(X/Q)\le t\le M$ corresponds to $1\le q\le Q$.  With
$r=\psi(t)=X/q$, for every measurable $|F_X|\le1$,
\begin{align*}
        \int_{\Phi(X/Q)}^M f(t)F_X(q(t))\dd t
        &=X\int_1^Q \frac{\Phi'(X/q)}{X/q}F_X(q)\frac{\dd q}{q^2}  \\
        &=\frac{2\lambda X}{Y}\int_1^Q F_X(q)\frac{\dd q}{q^2}
        +o_Q(X/Y).                                             \tag{18}
\end{align*}
Also, uniformly for $1\le q\le Q$,
\[
        G_X(\Phi(X/q))=q+o_Q(1),                               \tag{19}
\]
because $f$ varies by $o_Q(1)$ relatively on the interval of length $X$ ending
at $\Phi(X/q)$, while $Xf(\Phi(X/q))=q$ exactly.

Apply (18) with $F_X(q)=\fp{G_X(\Phi(X/q))}$.  By (19), the integrand converges
to $\fp q$ away from the finitely many integer discontinuities in $[1,Q]$;
small neighborhoods of those discontinuities have arbitrarily small
$q^{-2}\dd q$ mass.  Thus the saving is deterministic convergence under the
natural $q$-measure, not an equidistribution assertion.  The range $q>Q$ has
mass only $O(X/(QY))$, since
\[
        \int_X^{\Phi(X/Q)} f(t)\dd t
        =\int_{\psi(X)}^{X/Q}\frac{\Phi'(r)}r\dd r
        \ll \frac{X}{QY}.
\]
Letting $X\to\infty$ and then $Q\to\infty$,
\[
        E(X)=\frac{2\lambda X}{Y}
        \int_1^\infty\frac{\fp q}{q^2}\dd q+o(X/Y).
\]
Finally,
\[
        \int_1^\infty\frac{\fp q}{q^2}\dd q
        =\sum_{m\ge1}\int_m^{m+1}\frac{q-m}{q^2}\dd q
        =\sum_{m\ge1}\left(\log\frac{m+1}{m}-\frac1{m+1}\right)
        =1-\gamma,
\]
which proves (17).  On the $q$-scale the measure is
$(2\lambda X/Y)q^{-2}\dd q$, and the floor removes the average fractional part
$1-\gamma$.

\smallskip
\noindent\textit{6. Contradiction.}
Combining (10), (11), and (17),
\[
\begin{aligned}
        X
        &\le X+\lambda(3-2B-2\log(\log2))X/Y
        -2\lambda(1-\gamma)X/Y+o(X/Y) \\
        &=X-2\lambda(B-B_1)X/Y+o(X/Y).
\end{aligned}
\]
Since $B>B_1$, this is impossible for large $X$.  Thus (1) cannot hold
eventually.

For the endpoint limsup, apply the conclusion with any fixed $B^*>B_1$ along
an unbounded sequence $n_j$.  Then
\[
        \frac{a_{n_j}(\log n_j-\log\log n_j+B_1)}{n_j^2}
        >\lambda
        \frac{\log n_j-\log\log n_j+B_1}
             {\log n_j-\log\log n_j+B^*}
        \to\lambda,
\]
so the limsup is at least $\lambda$.
\end{proof}

\end{document}
